\begin{figure}[htbp]
 \centering
 \begin{tikzpicture}[scale=0.95,>=Latex]
  \fill[blue!6] (-2,0) circle (2);
  \draw[thick,blue!65!black] (-2,0) circle (2);
  \draw[->,gray!65] (-4.5,0)--(1.8,0);
  \draw[->,gray!65] (0,-2.35)--(0,2.5);
  \fill (-2,0) circle (1.5pt) node[below left] {$-RN(x_0)$};
  \fill (0,0) circle (1.7pt) node[below right] {$0\in T_{x_0}M$};
  \draw[thick,->] (0,0)--(1.25,0) node[above] {$N(x_0)$};
  \draw[thick] (-2,0)--(-2,2) node[midway,left] {$R$};
  \node at (-2,-1.1) {$U_0$};
  \node[align=center] at (-2,2.65) {A planar section of the tangent-space sphere};
 \end{tikzpicture}
 \caption{In dimension three, the open ball $U_0\subset T_{x_0}M$
 has radius $R$ and center $-RN(x_0)$, so the logarithm origin lies on
 its boundary. At this stage the perimeter support is known to lie on
 $\exp_{x_0}(\partial U_0)$. The phase-recovery argument identifies the occupied
 phase and then proves that the induced metric on the enclosed region
 is Euclidean.}
 \label{fig:offset-sphere}
\end{figure}
